
\chapter{Mastercard}\label{ch:mastercard}
От Visa Mastercard отличается единой сетевой архитектурой,
форматами GCMS (Global Clearing Management System)
и~IPM (Integrated Product Messages), токенизацией через
MDES (Mastercard Digital Enablement Service), механизмом
DSRP (Digital Secure Remote Payments) и~сервисом ABU
(Automatic Billing Updater)~--- при~общей четырёхсторонней
модели (см.~раздел~\ref{sec:ecosystem-four-party}), общем EMV
и~похожих таблицах interchange.

\importantnote{%
  \textbf{Контекст для~российского читателя.} С~марта 2022~года
  Mastercard приостановила обслуживание российских банков
  для~новых международных операций~\cite{mc-russia-suspension-2022}.
  Карты Mastercard, выпущенные российскими банками до~этого
  момента, продолжают приниматься внутри России. Авторизации
  и~клиринг по~таким картам внутри страны идут через
  Операционный платёжный и~клиринговый центр НСПК
  (ОПКЦ~НСПК; см.~гл.~\ref{ch:mir}), минуя Banknet~---
  глобальную сеть Mastercard для~авторизации, клиринга и~расчёта,
  а~сроки действия продлеваются российскими
  банками-эмитентами~\cite{cbr-russia-cards-normal-2022}.
  За~пределами РФ эти карты не~принимаются. Глава описывает
  архитектуру и~правила Mastercard как систему, чтобы объяснить
  логику, на~которую ориентирована унаследованная инфраструктура
  эквайеров в~России; местная специфика после 2022~года разбирается
  в~главе~\ref{ch:mir}.

}

\needspace{6\baselineskip}
\section{Сетевой периметр Mastercard}\label{sec:mc-banknet}

Сетевой домен Mastercard исторически называется Banknet
(см.~рис.~\ref{fig:banknet-topology})~--- именно через него идут
авторизация, клиринг и~расчёт между участниками. У~Mastercard
эти три фазы разнесены во~времени и~в~требованиях к~латентности
(см.~гл.~\ref{ch:tx-lifecycle}), и~каждую обслуживает своя подсистема.
\highlight{Switching Services описывают этот стек так: за~авторизацию отвечает
Authorization Platform, за~клиринг~--- Global Clearing Management
System, GCMS, за~расчёт~--- Settlement Account Management,
SAM}~\cite{mc-switching-services}.

Для~участника это означает один сетевой периметр, но~разные
интеграции: сертификация и~тестовые среды у~каждой подсистемы
свои~\cite{mc-rules}.

Недоступность Authorization Platform не~перенаправляет трафик
на~клиринг и~наоборот; недоступность авторизации обрабатывается
режимом \textit{stand-in} или приводит к~отказу, а~не к~переключению
на~альтернативную сеть~\cite{mc-switching-services}. Stand-In
Processing (STIP), не~отражённый на~схеме~\ref{fig:banknet-topology},
работает у~Mastercard по~той же логике, что и~у~Visa
(см.~раздел~\ref{sec:tx-authorization}): Authorization Platform отвечает
за~эмитента по~заранее согласованным
параметрам~\cite{mc-switching-services}.

На~периметре сети участники подключаются к~Authorization Platform
и~GCMS не~напрямую, а~через региональный edge-узел Mastercard
Interface Processor (MIP), обеспечивающий шифрование канала,
маршрутизацию и~базовую буферизацию обмена~\cite{mc-switching-services}.

\begin{figure}[htbp]
  \centering
  \includefiguresvg[width=\linewidth]{assets/figures/ch14-mastercard/banknet-topology}
  \caption{Топология логических подсистем сетевого периметра Mastercard.}\label{fig:banknet-topology}
\end{figure}
\FloatBarrier

\section{Dual message: IPM против TC33}\label{sec:mc-dmsa}

Общая логика dual message разобрана в~разделе~\ref{sec:tx-dual-single}.
\highlight{В~Mastercard авторизационный диалог строится вокруг
MTI~0100/MTI~0110 (MTI~--- message type identifier, четырёхзначный
код типа сообщения ISO~8583); на~клиринговом
этапе эквайер формирует для~GCMS интегрированные продуктовые
сообщения IPM~1240 (Integrated Product Messages). TC33
(формат пакетного клирингового файла Visa BASE~II), как у~Visa,
здесь не~применяется}~\cite{mc-chargeback-guide}. Физически IPM-файл
собирает эквайер на~основании авторизованных и~захваченных
операций; Authorization Platform сама IPM не~генерирует, она лишь
поставляет авторизационные ответы, на~которые опирается клиринг.

Наборы полей в~авторизации и~в~клиринге различаются.
DE25 (DE~--- data element, элемент данных ISO~8583 с~указанным номером
поля; см.~гл.~\ref{ch:iso8583}; DE25 несёт POS Condition Code)
расходится по~допустимым значениям, DE48
в~авторизации несёт один набор подэлементов, а~в~IPM расширяется
дополнительными subelement. DE25 и~DE48~--- лишь наиболее подвижные
примеры; полный набор расходящихся полей шире и~включает
DE2/3/4/14/22/43/48/49/61, а~также PDS-блоки (PDS~--- private data
subelement, приватные подэлементы данных Mastercard внутри
расширяемых полей IPM) внутри IPM. Если инженер при~отладке клирингового
отказа смотрит только в~авторизационный лог, реальная причина может
быть в~том, какие данные попали в~\textit{first presentment}~---
первое клиринговое предъявление операции эквайером в~адрес
эмитента~\cite{mc-chargeback-guide}.

\importantnote{%
  Dual message у~Visa и~Mastercard похож по~идее, но~инженерная единица
  клиринга разная: у~Visa BASE~II/TC33, у~Mastercard GCMS/IPM. Ошибка
  начинается, когда IPM трактуют как переименованный TC33; клиринговый
  конвейер Mastercard требует отдельной логики формирования и~разбора
  ответов.
}

\begin{figure}[htbp]
  \centering
  \includefiguresvg[width=.8\linewidth]{assets/figures/ch14-mastercard/auth-vs-gcms}
  \caption{Два потока Mastercard: Authorization Platform и~GCMS.}\label{fig:auth-vs-gcms}
\end{figure}
\FloatBarrier

\section{IPM и GCMS}\label{sec:mc-ipm}

GCMS принимает и~маршрутизирует данные клиринга в~формате IPM
(интегрированные продуктовые сообщения для~клиринга и~связанных
с~ним операционных процессов)~\cite{mc-switching-services}.

IPM~--- структурированный клиринговый формат со~своими типами
сообщений и~элементами данных~\cite{mc-chargeback-guide}. Подробная
карта записей, подэлементов и~служебных заголовков вынесена
в~специализированную документацию Mastercard: публичные материалы
сети подтверждают формат IPM и~роль GCMS, но~не дают полного
атласа полей для~построчного разбора.

\begin{table}[H]
  \begingroup
  \small
  \setlength{\tabcolsep}{4pt}
  \renewcommand{\arraystretch}{1.2}
  \begin{tabularx}{\textwidth}{
      @{}P{0.26\textwidth}
      Y@{}
    }
    \toprule
    \headrow
    \thd{Уровень} & \thd{Описание} \\
    \midrule
    Файл IPM & Контейнер клирингового обмена \\
    Клиринговая запись & Единица обмена внутри файла \\
    Message header & Служебный заголовок записи \\
    ISO~8583 data elements & Элементы данных транзакции \\
    Fees~/~adjustments & Комиссии и~корректировки \\
    \bottomrule
  \end{tabularx}%
  \endgroup
  \caption{Структура обмена IPM Mastercard.}\label{tab:mc-ipm-structure}
\end{table}

IPM~--- клиринговое сообщение Mastercard, отдельное от~авторизационного
ISO~8583. Его first presentment использует MTI~1240 вместо привычного
0100. IPM работает пакетно: получатель принимает файл целиком
и~обрабатывает содержимое асинхронно. Ответа уровня 1250 на~каждую
запись 1240 не~существует; подтверждение доставки даётся
административным сообщением MTI~1644 на~уровне файла, не~записи,
и~служит основой реконсиляции~\cite{mc-gcms-rm-2021}. 1240~---
\textit{directed}-сообщение без~смыслового ответа уровня клиринга:
подавляющее большинство first presentment проходит и~закрывается
расчётом нетто-позиции через SAM. IPM~1442 first chargeback~---
не~ответ на~1240, а~самостоятельная инициатива эмитента при~споре
клиента, со~своими основаниями и~сроками; 1442~--- тоже
\textit{directed}-сообщение, а~противодействие эквайера~--- 1240
second presentment с~функциональным кодом 205 (полная сумма) или~282
(часть)~\cite{mc-chargeback-guide,mc-tpr}. Набор элементов данных другой:
DE25 несёт POS Condition Code, DE48~--- расширенные данные
со~множеством подэлементов, включая токен, дополнительные коды
транзакции и~уже упомянутый Merchant Advice Code (DE48~SE84),
DE63~--- сетевые данные: маршрутизация и~флаги транзакции
внутри GCMS~\cite{mc-chargeback-guide}. Разработчик, переносящий
логику разбора с~авторизации на~клиринг, не~может переиспользовать
авторизационный парсер: структуры слишком разные.

До~марта 2022~г. Mastercard обслуживала внутрироссийские
транзакции; клиринговая запись для~покупки на~5\,000~руб. через чиповый
терминал выглядела так. В~клиринговом файле эквайер формировал
запись MTI~1240, \textit{first presentment}, заголовок которой содержал
Member Message Text и~процедурные поля, идентифицирующие участника
и~цикл обработки.
Далее шли data elements~--- DE~2 с~PAN, DE~4 с~суммой в~формате
\texttt{000000500000}, DE~49 с~кодом валюты \texttt{643},
DE~22 с~POS Entry Mode, указывающим на~чип, DE~48 с~подэлементами
SE~22 Multi-Purpose Merchant Indicator, SE~37 Additional Merchant
Data и~SE~84 Merchant Advice Code, если авторизация была
с~рекомендацией по~повторам. DE~63 нёс Network Data~---
флаги маршрутизации внутри GCMS. Отдельно появлялись
записи комиссий: interchange fee в~валюте расчёта
и~assessment fee сети~\cite{mc-chargeback-guide}.

Если эмитент оспаривает транзакцию, он создаёт \textit{first chargeback}
(первый возвратный платёж IPM~1442, инициатива эмитента в~адрес
эквайера), ссылающийся на~тот же DE~63 Network Data
и~Acquirer Reference Number (ARN~--- уникальный идентификатор
операции, присваиваемый эквайером для~сверки клиринга и~dispute).
Эквайер видит возвратный платёж в~следующем клиринговом файле;
для~\textit{second presentment} (повторное предъявление эквайером
после первого chargeback) он формирует IPM~1240 с~соответствующим
\textit{function code} 205 или~282 (коды цели сообщения IPM второго
представления) и~прикладывает доказательства через Mastercard
Connect~--- корпоративный портал Mastercard для~участников сети
с~доступом к~документации и~сервисам~\cite{mc-chargeback-guide}.

Код рекомендации для~ТСП (Merchant Advice Code, MAC)~--- рекомендация,
которую сеть передаёт в~авторизационном ответе в~DE~48, subelement~84
для~поддерживаемых сценариев~\cite{mc-tpr}. По~нему продавец или шлюз
решает, повторять ли попытку, обновлять ли реквизиты или прекращать
автоматические списания.

По~результатам клиринга сеть рассчитывает нетто-позиции участников
и~передаёт их в~расчётный процесс; конкретные сроки и~порядок
урегулирования задают правила региона и~роль участника~\cite{mc-rules}.
В~период между авторизацией и~расчётом сеть несёт расчётный риск:
правила Mastercard прямо гарантируют расчёт по~многим транзакциям
между клиентами, совокупный расчётный риск
(\textit{gross settlement exposure}) на~31~марта 2024~года
оценивается в~$73{,}8$~млрд~USD при~необеспеченном риске
$61{,}2$~млрд~\cite{mc-10k-settlement-risk-fy2024}. Участники
с~повышенным риском вносят денежное обеспечение, аккредитивы или
гарантии~\cite{mc-10k-settlement-risk-fy2024}; механизм по~функции
аналогичен Visa и~подробнее разобран в~разделе~\ref{sec:visa-settlement-guarantee}.

\section{Собственные реализации: M/Chip и MDES}\label{sec:mc-proprietary}

Поверх общих стандартов EMVCo Mastercard задаёт собственные
требования: M/Chip покрывает чиповое приложение на~карте и~терминале,
а~MDES~--- токенизацию и~удалённую оплату.

Общего знания EMV (гл.~\ref{ch:emv}) недостаточно, когда устройство,
карта или процессинг должны поддерживать Mastercard. Семейство M/Chip~\cite{mc-mchip}
поверх стандарта EMV добавляет собственные правила поведения приложения
и~отдельную процедуру сертификации. В~нём есть контактный и~бесконтактный
профили; M/Chip Contactless разобран в~разделе~\ref{sec:contactless-emv}.
Каждый профиль добавляет свой набор тегов, правил офлайн-поведения
и~тестовых сценариев.
Отсюда не~следует, что \enquote{если терминал уже поддерживает EMV,
то~Mastercard заработает автоматически}: нужны специфичные для~схемы
требования и~их проверка. Точную и~актуальную матрицу профилей
M/Chip сверяют через Mastercard Connect, в~открытом доступе полной
матрицы нет. Типичная ошибка при~сертификации: терминал проходит
EMVCo Level~2 (см.~гл.~\ref{ch:emv}), но~не проходит тесты M/Chip
из-за расхождения в~обработке конкретных тегов или~логике CVM
(cardholder verification method)~\cite{mc-mchip}.

Общая модель сетевого токена и~MDES как Token Service Provider
описана в~главе~\ref{ch:tokenization}. \highlight{Mastercard-специфика
в~MDES~--- встроенный механизм защищённой удалённой оплаты DSRP
(Digital Secure Remote Payments) для~приёма интернет-платежей
и~платежей в~приложениях с~опорой на~криптографию EMV}~\cite{mc-tpr}. DSRP формирует криптограмму
для~каждой операции CNP (card-not-present, без физического
предъявления карты) и~даёт эмитенту более сильное подтверждение
подлинности, чем обычный сценарий сохранённых
реквизитов~\cite{mc-mdes-token-connect}. Какой именно набор криптограмм
и~полей пришлёт клиентский кошелёк, зависит от~его реализации;
Mastercard публично не~фиксирует единый профиль для~всех Token
Requestors~--- зарегистрированных запрашивающих токены сторон
(кошельков, эквайеров, торговцев; см.~гл.~\ref{ch:tokenization}).

В~мобильных сценариях место выработки криптограммы задаёт модель
кошелька: при~Secure Element (SE) или Host Card Emulation с~локальными
ключами она формируется на~самом устройстве, при~облачном профиле~---
в~инфраструктуре TSP (Token Service Provider), а~устройство
получает однократный платёжный токен. С~точки зрения сети поле криптограммы выглядит одинаково,
но~операционные риски, требования к~аттестации устройства и~поведение
в~офлайне различаются.

M/Chip, MDES и~DSRP закрывают чиповый профиль и~токенизацию, но~не
исчерпывают CNP-периметр Mastercard: аутентификация держателя
при онлайн-оплате (3-D~Secure~2, см.~гл.~\ref{ch:3ds}) и~Click to~Pay
поверх EMVCo~SRC (Secure Remote Commerce~--- стандарт EMVCo
для~упрощённого онлайн-оформления заказа) на~схеме~\ref{fig:mchip-dsrp-layers}
не~отражены.

\begin{figure}[htbp]
  \centering
  \includefiguresvg[width=.7\linewidth]{assets/figures/ch14-mastercard/mchip-dsrp-layers}
  \caption{От EMV base к~M/Chip и~DSRP.}\label{fig:mchip-dsrp-layers}
\end{figure}
\FloatBarrier

\section{Mastercard Send}\label{sec:mc-send}

Mastercard Send~--- выплаты и~переводы, где средства
должны уйти получателю на~карту, банковский счёт, цифровой кошелёк
или другой канал, который поддерживает конкретная программа~\cite{mc-send-aptpay}.

Disbursements (выплаты физлицам от~бизнеса)~--- зарплаты, страховые
возмещения, выплаты продавцам маркетплейса и~платформам разовой
занятости (\textit{gig}).
P2P-переводы идут между физическими лицами как внутри страны,
так и~в~международных коридорах. В~карточном
сценарии платформа сначала инициирует финансирующую Funding
Transaction (резервирует средства отправителя), а~затем
запускает Payment Transaction~--- кредитовую операцию в~адрес
получателя~\cite{mc-tpr}.

Технически Mastercard Send требует от~участника поддержки
Send Transaction в~поле DE48. Авторизационный запрос на~выплату
несёт в~DE3 Transaction Processing Code, отличный от~обычной
покупки: транзакция не~дебетует карту, а~кредитует получателя.
Процессинговый парсер, настроенный только на~покупки и~возвраты,
транзакцию \textit{Send} не~обработает~--- нужна отдельная ветка логики.

Сравнительный аналог у~Visa~--- Visa Direct, такая же
push-модель с~той же парой Funding/Payment Transaction,
но~с~другими полями и~процедурами участника. Отличия видны
в~источнике финансирования и~в~нормах окончательности расчёта
(\textit{finality}): ряд коридоров Mastercard Send доставляет
выплаты почти в~реальном времени (\textit{near real-time}), за~секунды
или минуты, тогда как фактическая доступность средств у~получателя
зависит от~банка-получателя~\cite{mc-send-aptpay}.

Для~российского рынка после марта 2022~года Mastercard Send как
трансграничный (\textit{cross-border}) продукт для~участников из~РФ
операционно недоступен.
Внутренние операции по~картам Mastercard, выпущенным российскими
банками до~марта 2022~года, продолжают обслуживаться через НСПК
(см.~раздел~\ref{sec:mir-opkc}). Выплаты на~карту физлицу в~РФ 2026~года
обычно идут через B2C-переводы (от~бизнеса физлицу) Системы
быстрых платежей (СБП; см.~гл.~\ref{ch:sbp}), сбор средств между
своими счетами~--- через Me2Me Pull (запрос средств со~счёта
в~другом банке) и~Me2Me Push (отправка средств на~свой счёт
в~другом банке; см.~раздел~\ref{sec:sbp-scenarios}).

\practicenote{%
  Формулировка \textit{near real-time} в~маркетинге Mastercard Send
  не~равна фактической доступности средств получателю: сверяйте
  конкретный коридор и~SLA банка или провайдера-получателя.
}
\section{ABU: автоматическое обновление данных карты}\label{sec:mc-abu}

Общая модель обновления реквизитов после перевыпуска карты, обе
модели интеграции (batch и~real-time API), сетевые токены как альтернатива
и~аналог Visa Account Updater (VAU) разобраны
в~разделе~\ref{sec:subscriptions-account-updater}. Mastercard-специфика
Automatic Billing Updater (ABU)~--- в~ограничениях охвата.
Сервис покрывает не~все типы карт; предоплаченные и~часть
региональных продуктов могут оставаться вне обновления; участие
эмитента в~ABU обязательным не~является~\cite{mc-abu,mc-tpr}.
Технические нормативные требования Mastercard (Technical
Product Requirements, TPR), на~которые ссылается эта и~следующие
секции, доступны участникам сети через Mastercard
Connect~\cite{mc-tpr}.

Платформа, работающая и~в~Visa, и~в~Mastercard, поддерживает оба
сервиса отдельно: операционные процедуры ABU и~VAU не~совпадают.

Структурная альтернатива ABU~--- сетевые токены MDES: под токеном
PAN устройства/реквизитов остаётся стабильным, а~смена базового PAN
происходит на~стороне TSP, поэтому отдельный механизм рассылки
обновлений эквайерам и~торговцам не~требуется. Для~рекуррентных
сценариев с~сохранёнными токенами ABU теряет смысл.

\section{Interchange Mastercard}\label{sec:mc-interchange}

Общий каркас interchange (регион~$\to$~продукт~$\to$~MCC~$\to$
канал~$\to$~квалифицирующие условия~$\to$~downgrade), MSC и~тарифные
модели для~продавца разобраны в~разделе~\ref{sec:ecosystem-interchange}
и~применимы к~Mastercard без изменений~\cite{mc-interchange-us}.
Сетевая специфика Mastercard~--- в~ярлыке \emph{Data Quality},
лестничной механике downgrade и~обновлении правил
по~magstripe-fallback.

Признаки, по~которым сеть квалифицирует транзакцию, удобно группировать
под общим ярлыком \emph{Data Quality}, объединяющим классы проверок:
ECI (electronic commerce indicator)/CVV2 (Card Verification Value~2)/
AVS (Address Verification Service) для~CNP (см.~гл.~\ref{ch:3ds}),
тайминг авторизации и~клиринга, POS Entry Mode и~CVM-исход
(cardholder verification method) для~card-present, признаки
card-on-file (сохранённые у~продавца реквизиты для повторных
списаний; см.~гл.~\ref{ch:tokenization}) и~сетевого токена.
Нарушение в~любом классе может сместить транзакцию
в~менее выгодную категорию~\cite{mc-interchange-us}.

Downgrade возникает, когда транзакция не~набирает достаточно
квалифицирующих признаков для~целевой категории. Типичные причины:
отсутствие electronic commerce indicator (ECI), неполные данные
в~DE48, пропущенный Address Verification Service (AVS) или
позднее предъявление в~клиринг (\textit{late presentment}~---
клиринговое предъявление операции с~опозданием против сроков,
заданных правилами схемы). Downgrade применяется ступенчато:
каждый пропущенный признак опускает транзакцию на~одну ступень
по~шкале ставок, в~предельном случае~--- в~Standard Rate,
самую дорогую из~общих категорий~\cite{mc-interchange-us}.

Для~аудита downgrade нужен доступ к~клиринговым данным: в~файле IPM
видно, какую категорию присвоила сеть и~какие признаки отсутствовали.
Авторизационный лог для~этого недостаточен, квалификация происходит
на~этапе клиринга~\cite{mc-interchange-us}.

В~апрельском релизе правил Mastercard 2026~года анонсирован
дополнительный сбор Fallback Avoidance~--- 0{,}10\,\% от~суммы
за~избыточное использование magstripe-fallback там, где доступен чип
или contactless, в~США по~Mastercard core (основной
кредитно-дебетовый бренд) и~Maestro (исторический дебетовый
бренд Mastercard). Публичные обзоры расходятся по~дате pass-through
к~продавцу (1~января против 1~апреля 2026~года): Mastercard ввела
правило на~уровне сети с~января, но~часть эквайеров переносит сбор
в~ставку с~апреля~\cite{mc-merchant-rates-2025-2026}.

\begin{figure}[tbp]
  \centering
  \includefiguresvg[width=.8\linewidth]{assets/figures/ch14-mastercard/interchange-qualification}
  \caption{Дерево квалификации interchange: от~региона до~даунгрейда.}\label{fig:interchange-qualification}
\end{figure}

Без привязки к~конкретному рынку, дате публикации и~публичной
региональной таблице нельзя надёжно утверждать, что тот или иной
премиальный продукт Mastercard \emph{всегда} дороже сопоставимого
продукта Visa. Такие сравнения быстро превращаются в~смесь локальных
правил, частных договорённостей и~устаревших тарифных редакций.
Для~рынка США Mastercard публикует отдельные таблицы ставок
и~критериев interchange; любой численный пример привязывается
к~конкретной региональной таблице~\cite{mc-interchange-us}.

\section{Архитектурные отличия от~Visa}\label{sec:mc-vs-visa}

Различия между Mastercard и~Visa бывают трёх родов: терминологические, исторические
и~инженерные. Терминологические и~исторические достаточно знать;
инженерные определяют, какой код придётся переписать при~переносе
системы с~одной схемы на~другую.

\begin{table}[H]
  \begingroup
  \small
  \setlength{\tabcolsep}{4pt}
  \renewcommand{\arraystretch}{1.2}
  \begin{tabularx}{\textwidth}{
      @{}P{0.26\textwidth}
      Y
      Y@{}
    }
    \toprule
    \headrow
    \thd{Параметр} & \thd{Mastercard}
    & \thd{Visa} \\
    \midrule
    Сетевая архитектура & Единый сетевой периметр Mastercard
    & BASE~I для~авторизации и~BASE~II для~клиринга \\
    Клиринг & IPM через~GCMS
    & TC33 через~BASE~II \\
    Модель dual message & Авторизация и~клиринг разнесены внутри
    одной сети & Авторизация и~клиринг разнесены по~двум именованным сетям \\
    Токенизация & MDES~+~DSRP
    & VTS (Visa Token Service)~+~Token Assurance Level \\
    Чиповые требования & M/Chip
    & Visa ICC~/~VCPS (Visa Contactless Payment Specification) \\
    Выплаты & Mastercard Send
    & Visa Direct \\
    Обновление реквизитов & Automatic Billing Updater (ABU)
    & Visa Account Updater (VAU) \\
    \bottomrule
  \end{tabularx}%
  \endgroup
  \caption{Архитектурные отличия Mastercard и~Visa.}\label{tab:mc-vs-visa}
\end{table}

Существеннее всего различия в~клиринговой модели и~формате обмена.
Разработчик, переносящий систему обработки клиринга с~одной схемы
на~другую, не~может переименовать поля: структуры IPM и~TC33
различаются настолько, что требуют отдельного парсинга и~отдельной
операционной диагностики. Различия в~токенизации важны
для~разработчиков кошельков и~интернет-платёжных интеграций. Различия
в~ABU и~VAU важны для~платёжных платформ с~рекуррентными платежами.

Есть и~различия другого рода: одни и~те же отраслевые слова
у~схем означают разное. Термин \enquote{dual message} применяют
и~к~Visa, и~к~Mastercard, но~за~ним стоят разные технические
реализации; общая формула \enquote{авторизация и~клиринг разделены
во~времени} верна для~обеих схем, расходятся детали разделения.

\enlargethispage{4\baselineskip}
\practicenote{%
  При работе с~сообщениями IPM Mastercard для~диагностики отказов
  нужно смотреть на~сам факт отказа и~на~Merchant Advice Code
  в~авторизационном ответе: именно он подсказывает, допустимы
  ли повторы и~нужно ли обновлять сохранённые реквизиты~\cite{mc-tpr}.

}
